% ======================================================================
%  AT-MONO113 — Chapter 13: Boundaries of The Actualization Theory
%  Canonical Monograph (v2.0) · The Actualization Theory:
%  A Reconstruction of Physics from Difference, Actualization and Spectrum
%
%  Status: COMPLETE — PASS-READY (MONO007 referee-ready architecture)
%  Inputs: Chapter01_Difference.tex ... Chapter06_D96Spectrum.tex,
%          ATQG_CanonicalMonograph.md, ATQG_Mono007RefereeReadiness.md,
%          ATVALID_UntestedItemsAudit.md
%  Method: assembly only — canonical results only
%  Citations: QG185 (Bekenstein quarter), QG196 (impossibility proof),
%             QG242 (SM dynamics hosted), QG245 (SM dynamics not complete),
%             QG278 (fundamental boundary), QG286 (difference duality),
%             QG291 (framework necessity), QG299 (post-frontier),
%             MONO007 (referee readiness), VALID001 (validation classification)
%
%  Usage: % ======================================================================
%  AT-MONO113 — Chapter 13: Boundaries of The Actualization Theory
%  Canonical Monograph (v2.0) · The Actualization Theory:
%  A Reconstruction of Physics from Difference, Actualization and Spectrum
%
%  Status: COMPLETE — PASS-READY (MONO007 referee-ready architecture)
%  Inputs: Chapter01_Difference.tex ... Chapter06_D96Spectrum.tex,
%          ATQG_CanonicalMonograph.md, ATQG_Mono007RefereeReadiness.md,
%          ATVALID_UntestedItemsAudit.md
%  Method: assembly only — canonical results only
%  Citations: QG185 (Bekenstein quarter), QG196 (impossibility proof),
%             QG242 (SM dynamics hosted), QG245 (SM dynamics not complete),
%             QG278 (fundamental boundary), QG286 (difference duality),
%             QG291 (framework necessity), QG299 (post-frontier),
%             MONO007 (referee readiness), VALID001 (validation classification)
%
%  Usage: % ======================================================================
%  AT-MONO113 — Chapter 13: Boundaries of The Actualization Theory
%  Canonical Monograph (v2.0) · The Actualization Theory:
%  A Reconstruction of Physics from Difference, Actualization and Spectrum
%
%  Status: COMPLETE — PASS-READY (MONO007 referee-ready architecture)
%  Inputs: Chapter01_Difference.tex ... Chapter06_D96Spectrum.tex,
%          ATQG_CanonicalMonograph.md, ATQG_Mono007RefereeReadiness.md,
%          ATVALID_UntestedItemsAudit.md
%  Method: assembly only — canonical results only
%  Citations: QG185 (Bekenstein quarter), QG196 (impossibility proof),
%             QG242 (SM dynamics hosted), QG245 (SM dynamics not complete),
%             QG278 (fundamental boundary), QG286 (difference duality),
%             QG291 (framework necessity), QG299 (post-frontier),
%             MONO007 (referee readiness), VALID001 (validation classification)
%
%  Usage: % ======================================================================
%  AT-MONO113 — Chapter 13: Boundaries of The Actualization Theory
%  Canonical Monograph (v2.0) · The Actualization Theory:
%  A Reconstruction of Physics from Difference, Actualization and Spectrum
%
%  Status: COMPLETE — PASS-READY (MONO007 referee-ready architecture)
%  Inputs: Chapter01_Difference.tex ... Chapter06_D96Spectrum.tex,
%          ATQG_CanonicalMonograph.md, ATQG_Mono007RefereeReadiness.md,
%          ATVALID_UntestedItemsAudit.md
%  Method: assembly only — canonical results only
%  Citations: QG185 (Bekenstein quarter), QG196 (impossibility proof),
%             QG242 (SM dynamics hosted), QG245 (SM dynamics not complete),
%             QG278 (fundamental boundary), QG286 (difference duality),
%             QG291 (framework necessity), QG299 (post-frontier),
%             MONO007 (referee readiness), VALID001 (validation classification)
%
%  Usage: \input{Chapter13_BoundariesOfTheActualizationTheory.tex} after the
%         preamble that defines the theorem environments (or include via the
%         master file).
% ======================================================================

% ---- Theorem environments (define in the master preamble if not present) ----
% \newtheorem{definition}{Definition}[chapter]
% \newtheorem{lemma}[definition]{Lemma}
% \newtheorem{theorem}[definition]{Theorem}
% \newtheorem{corollary}[definition]{Corollary}
% \newtheorem{remark}[definition]{Remark}

\chapter{Boundaries of The Actualization Theory}
\label{c13:chap:boundaries}

\section{Introduction}
\label{c13:sec:boundaries-introduction}

The preceding chapters derived the physics of The Actualization Theory: the foundation, the
spectrum, quantum mechanics, gravity, the standard model, and cosmology. This chapter is the
Boundary Layer of the monograph: it separates \emph{ontological boundaries} from \emph{physics
derivation}, documents the limits of the framework, and shows that the boundaries are
\emph{not derivation failures} --- they are the explicit limits beyond which the framework
does not claim to derive.

The boundaries documented here are: Difference as an ontological boundary \cite{c13:QG278}; the
status of the tensor face $\psi$ \cite{c13:QG286}; the numerical boundary $\pi$ \cite{c13:QG291};
the Bekenstein $2\pi$ boundary \cite{c13:QG185,c13:QG196}; and the hosted standard-model dynamics
\cite{c13:QG242,c13:QG245}. The chapter draws the distinction between \emph{derived physics} and
\emph{boundary questions} \cite{c13:QG299} and states explicitly what the theory does not claim.

Following VALID001 \cite{c13:VALID001}, \emph{experimental validation is not a boundary}: the
derived predictions $S,T,U$, the electron g-2 $a_e$, and the Majorana character $0\nu\beta\beta$
await experiment and are not listed as derivation gaps. Consistent with MONO007 \cite{c13:MONO007},
the presentation is referee-safe and introduces no new physics and no speculation.

\section{The Nature of Scientific Boundaries}
\label{c13:sec:nature-boundaries}

\begin{lemma}[Scientific boundaries are not failures]
\label{c13:lem:boundaries-not-failures}
A \emph{scientific boundary} is a limit of a framework beyond which the framework does not
claim to derive: either an irreducible primitive (an ontological boundary), a numerical
constant not derivable inside the framework (a numerical boundary), or a hosted structure
(an external component carried without derivation). A documented boundary is not a
derivation failure: it is the explicit, honest limit of the framework, disclosed so that no
claim exceeds what is established.
\end{lemma}

\begin{proof}
A framework is characterized as much by what it does not claim as by what it derives: the
disclosure of a boundary is the statement "this framework does not derive X". The
post-frontier audit \cite{c13:QG299} classifies the remaining items as open, partial,
boundary, or methodology, distinguishing the boundaries from the open derivation items.
Hence a boundary is an honest limit, not a failure of the derivation.
\end{proof}

\section{Difference as an Ontological Boundary}
\label{c13:sec:difference-boundary}

Difference is the ontological boundary of the theory: the irreducible primitive whose
count-conservation law is its definitional identity, and which no deeper notion explains
--- the \emph{primitive} boundary, the irreducible starting point rather than an unexplained
gap. The fundamental-boundary audit \cite{c13:QG278} establishes that every attempt to reduce
it reintroduces it (two distinct things use distinct = difference; a boundary is where
things differ; a transition is a before/after difference); by
Chapter~\ref{c01:chap:difference} (Theorem~\ref{c01:thm:fundamental-boundary}) it is the
fundamental boundary of the theory, and by Lemma~\ref{c13:lem:boundaries-not-failures} an
ontological boundary is an honest primitive limit, not a derivation failure.
\label{c13:thm:difference-ontological}

\section{The Status of \texorpdfstring{$\psi$}{psi}}
\label{c13:sec:psi-status}

\begin{theorem}[\texorpdfstring{$\psi$}{psi} is the tensor face of Difference]
\label{c13:thm:psi-face}
The tensor field $\psi$ is the tensor face of Difference: the traceless part of the
trace/traceless decomposition of the one Difference, read against the tensor reference
$\eta$. Its deepest ontological nature --- whether the tensor content is ultimately
irreducible to the scalar count --- is unresolved: a boundary question, not a derivation gap.
\end{theorem}

\begin{proof}
The Difference duality \cite{c13:QG286} establishes that $\rho$ and $\psi$ are the trace and
traceless faces of the one Difference (Chapter~\ref{c01:chap:difference},
Theorem~\ref{c01:thm:duality}); the tensor face is read against the reference $\eta$
(Chapter~\ref{c02:chap:eta}, Theorem~\ref{c02:thm:tensor-read}). Hence $\psi$ is the tensor
face of Difference --- not an independent primitive. The post-frontier audit
\cite{c13:QG299} classifies the $\psi$ question as a boundary item: the observables of the
tensor face are derived, but its ultimate ontological status is a boundary question.
\end{proof}

\section{The \texorpdfstring{$\pi$}{pi} Boundary}
\label{c13:sec:pi-boundary}

\begin{theorem}[\texorpdfstring{$\pi$}{pi} is redundant as a theory input]
\label{c13:thm:pi-redundant}
\emph{$\pi$} is the universal mathematical constant, the geometric ratio of the circle; in
the canonical architecture it is a \emph{numerical boundary} --- a constant not derivable
inside the framework --- and not a primitive. Correspondingly, $\pi$ is redundant as a
theory input: no derived prediction uses $\pi$ as an input, and where it appears it is an
imported numerical factor.
\end{theorem}

\begin{proof}
The framework-necessity audit \cite{c13:QG291} establishes the classification: every derived
observable --- masses, couplings, mixings, cosmological fractions, spectral indices,
acoustic ratios, and the pre-registered predictions --- is a pure ratio of D96 spectral
constants together with the calibration scale. $\pi$ appears only where an imported quantum
factor is required, most notably in the Bekenstein coefficient; as a constant not derivable
inside the framework it is a numerical boundary, distinct from the primitive reference
$\eta$ (Chapter~\ref{c02:chap:eta}, Theorem~\ref{c02:thm:eta-primitive-pi-boundary}), not a
primitive.
\end{proof}

\section{The Bekenstein \texorpdfstring{$2\pi$}{2pi} Boundary}
\label{c13:sec:bekenstein}

\begin{theorem}[The Bekenstein $1/4$ coefficient requires the imported $2\pi$ factor]
\label{c13:thm:bekenstein-2pi}
The exact Bekenstein coefficient $S = A/4$ is a boundary: the structure $S \propto A$,
$M \propto R$, $T \propto 1/R$ is derived, but the exact $1/4$ coefficient requires the
imported quantum factor $T = \kappa/(2\pi)$, which the framework cannot produce internally.
The boundary is therefore explicitly referenced to the imported quantum factor, not
presented as a derivation of the theory.
\end{theorem}

\begin{proof}
The Bekenstein origin \cite{c13:QG185} derives the structure ($S \propto A$, $M \propto R$,
$T \propto 1/R$); the deficit first-law gives $S = A/(8\pi)$, not $1/4$, and the exact $1/4$
requires the $2\pi$ quantum factor $T = \kappa/(2\pi)$, absent in the framework. The
impossibility proof \cite{c13:QG196} confirms that the exact $1/4$ cannot be derived without
importing $\pi$. Consistent with Chapter~\ref{c10:chap:gravity}
(Theorem~\ref{c10:thm:bekenstein-boundary}), the boundary is referenced, not claimed.
\end{proof}

\section{Hosted Standard Model Dynamics}
\label{c13:sec:hosted-sm}

\begin{theorem}[The SM Lagrangian is hosted, not derived]
\label{c13:thm:sm-hosted}
The standard-model \emph{Lagrangian} and its dynamics are hosted: the framework derives the
observables (masses, couplings, mixings) but carries the Lagrangian form as an external
structure. The hosted SM dynamics (the Lagrangian and its structure) is thus a boundary,
while the derived SM observables (the values read from the D96 moments) are derived physics.
\end{theorem}

\begin{proof}
The SM-dynamics audits \cite{c13:QG242,c13:QG245} establish the classification: the gauge
\emph{symmetry} is derived, but the gauge \emph{dynamics} (the Lagrangian, vertices,
propagators) is hosted; the SM dynamics is not complete as a derivation. The derived
observables --- the fermion sectors, couplings, mixings, weak scale, and precision
predictions (Chapter~\ref{c11:chap:sm}) --- are reads of the D96 moments. Hence the SM
dynamics is a boundary and the observables are derived physics.
\end{proof}

\section{Derived Physics vs Boundary Questions}
\label{c13:sec:derived-vs-boundary}

\begin{theorem}[The classification of remaining items]
\label{c13:thm:classification}
The remaining items of the framework classify into derived physics, experimental
validation, and boundaries: the derived physics is established; the validation items
($S,T,U$, $a_e$, $0\nu\beta\beta$) are derived predictions awaiting experiment; the
boundaries are the ontological, numerical, and hosted limits.
\end{theorem}

\begin{proof}
The post-frontier audit \cite{c13:QG299} classifies the remaining items; VALID001
\cite{c13:VALID001} establishes that $S,T,U$, $a_e$, and $0\nu\beta\beta$ are \emph{derived
predictions} (status A) belonging to the experimental-validation category, with no missing
derivation steps. Hence none is an open derivation issue or a boundary; the boundaries are
those documented in this chapter, and the remaining items partition into derived physics
(established), validation (derived, awaiting experiment), and boundaries (the limits).
\end{proof}

\begin{corollary}[No validation item is an open derivation issue]
\label{c13:cor:no-open-validation}
The quantities $S,T,U$, the electron g-2 $a_e$, and the Majorana character $0\nu\beta\beta$
are not listed as open derivation issues: their derivations are complete, and their
experimental validation is the open task.
\end{corollary}

\section{What Is Not Claimed}
\label{c13:sec:not-claimed}

\begin{theorem}[The framework's non-claims]
\label{c13:thm:non-claims}
The framework does not claim: (i) a derivation of the standard-model Lagrangian (hosted);
(ii) a derivation of the exact Bekenstein $1/4$ coefficient (requires the imported $2\pi$
factor); (iii) a resolution of the ultimate ontological status of $\psi$; (iv) any result
beyond the accepted derivations. Equivalently, the framework claims exactly what it
derives: the physics observables established in the preceding chapters, with the boundaries
disclosed and the validation items awaiting experiment.
\end{theorem}

\begin{proof}
(i) The SM dynamics is hosted \cite{c13:QG242,c13:QG245} (Theorem~\ref{c13:thm:sm-hosted}). (ii) The
Bekenstein $1/4$ coefficient requires the imported $2\pi$ factor \cite{c13:QG185,c13:QG196}
(Theorem~\ref{c13:thm:bekenstein-2pi}). (iii) The ontological status of $\psi$ is unresolved
(Theorem~\ref{c13:thm:psi-face}). (iv) The monograph is assembly-only, restricted to
accepted results. Hence the non-claims are explicit and the claimed content is the derived
physics.
\end{proof}

\section{Consequences for Future Work}
\label{c13:sec:future-work}

\begin{theorem}[The boundaries delimit the future work]
\label{c13:thm:future-work}
The boundaries delimit the future work of the framework: the experimental validation of the
derived predictions, the resolution of the ontological status of $\psi$, and the question
of whether the imported $2\pi$ factor can be internalized. The boundaries are not blockers:
the derived content of the theory is complete within its scope, and the boundaries mark
where the framework honestly stops.
\end{theorem}

\begin{proof}
The derived predictions ($S,T,U$, $a_e$, $0\nu\beta\beta$) await experimental validation
\cite{c13:VALID001}; the ontological status of $\psi$ is a boundary question
(Theorem~\ref{c13:thm:psi-face}); the Bekenstein $2\pi$ factor is an imported boundary
(Theorem~\ref{c13:thm:bekenstein-2pi}). By Lemma~\ref{c13:lem:boundaries-not-failures} the
boundaries are not failures, and the derived physics is complete within the canonical
hierarchy (Chapter~\ref{c12:chap:cosmology}, Corollary~\ref{c12:cor:canonical-cosmology}).
Hence the boundaries are not blockers but the framework's declared edge.
\end{proof}

\section{Summary}
\label{c13:sec:boundaries-summary}

This chapter has documented the boundaries of the theory and shown that they are honest
limits, not derivation failures. Difference is the irreducible primitive boundary; $\psi$ is
the tensor face of Difference with unresolved ontological status; $\pi$ is a numerical
boundary, redundant as a theory input; the Bekenstein $1/4$ coefficient is explicitly tied
to the imported $2\pi$ factor; and the SM Lagrangian is hosted while the observables are
derived. The remaining items partition into derived physics, validation, and boundary, with
the validation items ($S,T,U$, $a_e$, $0\nu\beta\beta$) distinct from open derivation issues
(Theorem~\ref{c13:thm:classification}, Corollary~\ref{c13:cor:no-open-validation}); the
framework claims exactly its derivations (Theorem~\ref{c13:thm:non-claims}); and the
boundaries delimit the honest future work without blocking it
(Theorem~\ref{c13:thm:future-work}).

The boundaries of The Actualization Theory are therefore: Difference (the ontological
boundary and primitive), $\psi$ (the tensor face with unresolved ontological status),
$\pi$ (the numerical boundary), the Bekenstein $2\pi$ factor (the imported quantum
boundary), and the hosted standard-model dynamics --- each an explicit limit, none a
derivation failure. The derived physics is complete within its scope; the experimental
validation of the derived predictions is the open task; and the boundaries mark where the
framework honestly stops. This closes the Boundary Layer; the following chapter documents
the frontier and the falsification paths.

\begin{thebibliography}{99}
\bibitem{c13:QG185} AT-QG Phase 185, \emph{Bekenstein Quarter Coefficient Origin}, Docs/Research.
\bibitem{c13:QG196} AT-QG Phase 196, \emph{Quarter Coefficient Impossibility Proof}, Docs/Research.
\bibitem{c13:QG242} AT-QG Phase 242, \emph{Standard Model Dynamics Audit}, Docs/Research.
\bibitem{c13:QG245} AT-QG Phase 245, \emph{Higgs Yukawa Origin Audit}, Docs/Research.
\bibitem{c13:QG278} AT-QG Phase 278, \emph{Fundamental Boundary Audit}, Docs/Research.
\bibitem{c13:QG286} AT-QG Phase 286, \emph{Difference Duality Audit}, Docs/Research.
\bibitem{c13:QG291} AT-QG Phase 291, \emph{Framework Necessity Audit}, Docs/Research.
\bibitem{c13:QG299} AT-QG Phase 299, \emph{Remaining Frontier Re-audit}, Docs/Research.
\bibitem{c13:MONO007} AT-MONO007, \emph{Referee-Readiness Resolution}, Docs/Zenodo/Monograph.
\bibitem{c13:VALID001} AT-VALID001, \emph{Remaining Untested Items Audit}, Docs/Research.
\end{thebibliography}
 after the
%         preamble that defines the theorem environments (or include via the
%         master file).
% ======================================================================

% ---- Theorem environments (define in the master preamble if not present) ----
% \newtheorem{definition}{Definition}[chapter]
% \newtheorem{lemma}[definition]{Lemma}
% \newtheorem{theorem}[definition]{Theorem}
% \newtheorem{corollary}[definition]{Corollary}
% \newtheorem{remark}[definition]{Remark}

\chapter{Boundaries of The Actualization Theory}
\label{c13:chap:boundaries}

\section{Introduction}
\label{c13:sec:boundaries-introduction}

The preceding chapters derived the physics of The Actualization Theory: the foundation, the
spectrum, quantum mechanics, gravity, the standard model, and cosmology. This chapter is the
Boundary Layer of the monograph: it separates \emph{ontological boundaries} from \emph{physics
derivation}, documents the limits of the framework, and shows that the boundaries are
\emph{not derivation failures} --- they are the explicit limits beyond which the framework
does not claim to derive.

The boundaries documented here are: Difference as an ontological boundary \cite{c13:QG278}; the
status of the tensor face $\psi$ \cite{c13:QG286}; the numerical boundary $\pi$ \cite{c13:QG291};
the Bekenstein $2\pi$ boundary \cite{c13:QG185,c13:QG196}; and the hosted standard-model dynamics
\cite{c13:QG242,c13:QG245}. The chapter draws the distinction between \emph{derived physics} and
\emph{boundary questions} \cite{c13:QG299} and states explicitly what the theory does not claim.

Following VALID001 \cite{c13:VALID001}, \emph{experimental validation is not a boundary}: the
derived predictions $S,T,U$, the electron g-2 $a_e$, and the Majorana character $0\nu\beta\beta$
await experiment and are not listed as derivation gaps. Consistent with MONO007 \cite{c13:MONO007},
the presentation is referee-safe and introduces no new physics and no speculation.

\section{The Nature of Scientific Boundaries}
\label{c13:sec:nature-boundaries}

\begin{lemma}[Scientific boundaries are not failures]
\label{c13:lem:boundaries-not-failures}
A \emph{scientific boundary} is a limit of a framework beyond which the framework does not
claim to derive: either an irreducible primitive (an ontological boundary), a numerical
constant not derivable inside the framework (a numerical boundary), or a hosted structure
(an external component carried without derivation). A documented boundary is not a
derivation failure: it is the explicit, honest limit of the framework, disclosed so that no
claim exceeds what is established.
\end{lemma}

\begin{proof}
A framework is characterized as much by what it does not claim as by what it derives: the
disclosure of a boundary is the statement "this framework does not derive X". The
post-frontier audit \cite{c13:QG299} classifies the remaining items as open, partial,
boundary, or methodology, distinguishing the boundaries from the open derivation items.
Hence a boundary is an honest limit, not a failure of the derivation.
\end{proof}

\section{Difference as an Ontological Boundary}
\label{c13:sec:difference-boundary}

Difference is the ontological boundary of the theory: the irreducible primitive whose
count-conservation law is its definitional identity, and which no deeper notion explains
--- the \emph{primitive} boundary, the irreducible starting point rather than an unexplained
gap. The fundamental-boundary audit \cite{c13:QG278} establishes that every attempt to reduce
it reintroduces it (two distinct things use distinct = difference; a boundary is where
things differ; a transition is a before/after difference); by
Chapter~\ref{c01:chap:difference} (Theorem~\ref{c01:thm:fundamental-boundary}) it is the
fundamental boundary of the theory, and by Lemma~\ref{c13:lem:boundaries-not-failures} an
ontological boundary is an honest primitive limit, not a derivation failure.
\label{c13:thm:difference-ontological}

\section{The Status of \texorpdfstring{$\psi$}{psi}}
\label{c13:sec:psi-status}

\begin{theorem}[\texorpdfstring{$\psi$}{psi} is the tensor face of Difference]
\label{c13:thm:psi-face}
The tensor field $\psi$ is the tensor face of Difference: the traceless part of the
trace/traceless decomposition of the one Difference, read against the tensor reference
$\eta$. Its deepest ontological nature --- whether the tensor content is ultimately
irreducible to the scalar count --- is unresolved: a boundary question, not a derivation gap.
\end{theorem}

\begin{proof}
The Difference duality \cite{c13:QG286} establishes that $\rho$ and $\psi$ are the trace and
traceless faces of the one Difference (Chapter~\ref{c01:chap:difference},
Theorem~\ref{c01:thm:duality}); the tensor face is read against the reference $\eta$
(Chapter~\ref{c02:chap:eta}, Theorem~\ref{c02:thm:tensor-read}). Hence $\psi$ is the tensor
face of Difference --- not an independent primitive. The post-frontier audit
\cite{c13:QG299} classifies the $\psi$ question as a boundary item: the observables of the
tensor face are derived, but its ultimate ontological status is a boundary question.
\end{proof}

\section{The \texorpdfstring{$\pi$}{pi} Boundary}
\label{c13:sec:pi-boundary}

\begin{theorem}[\texorpdfstring{$\pi$}{pi} is redundant as a theory input]
\label{c13:thm:pi-redundant}
\emph{$\pi$} is the universal mathematical constant, the geometric ratio of the circle; in
the canonical architecture it is a \emph{numerical boundary} --- a constant not derivable
inside the framework --- and not a primitive. Correspondingly, $\pi$ is redundant as a
theory input: no derived prediction uses $\pi$ as an input, and where it appears it is an
imported numerical factor.
\end{theorem}

\begin{proof}
The framework-necessity audit \cite{c13:QG291} establishes the classification: every derived
observable --- masses, couplings, mixings, cosmological fractions, spectral indices,
acoustic ratios, and the pre-registered predictions --- is a pure ratio of D96 spectral
constants together with the calibration scale. $\pi$ appears only where an imported quantum
factor is required, most notably in the Bekenstein coefficient; as a constant not derivable
inside the framework it is a numerical boundary, distinct from the primitive reference
$\eta$ (Chapter~\ref{c02:chap:eta}, Theorem~\ref{c02:thm:eta-primitive-pi-boundary}), not a
primitive.
\end{proof}

\section{The Bekenstein \texorpdfstring{$2\pi$}{2pi} Boundary}
\label{c13:sec:bekenstein}

\begin{theorem}[The Bekenstein $1/4$ coefficient requires the imported $2\pi$ factor]
\label{c13:thm:bekenstein-2pi}
The exact Bekenstein coefficient $S = A/4$ is a boundary: the structure $S \propto A$,
$M \propto R$, $T \propto 1/R$ is derived, but the exact $1/4$ coefficient requires the
imported quantum factor $T = \kappa/(2\pi)$, which the framework cannot produce internally.
The boundary is therefore explicitly referenced to the imported quantum factor, not
presented as a derivation of the theory.
\end{theorem}

\begin{proof}
The Bekenstein origin \cite{c13:QG185} derives the structure ($S \propto A$, $M \propto R$,
$T \propto 1/R$); the deficit first-law gives $S = A/(8\pi)$, not $1/4$, and the exact $1/4$
requires the $2\pi$ quantum factor $T = \kappa/(2\pi)$, absent in the framework. The
impossibility proof \cite{c13:QG196} confirms that the exact $1/4$ cannot be derived without
importing $\pi$. Consistent with Chapter~\ref{c10:chap:gravity}
(Theorem~\ref{c10:thm:bekenstein-boundary}), the boundary is referenced, not claimed.
\end{proof}

\section{Hosted Standard Model Dynamics}
\label{c13:sec:hosted-sm}

\begin{theorem}[The SM Lagrangian is hosted, not derived]
\label{c13:thm:sm-hosted}
The standard-model \emph{Lagrangian} and its dynamics are hosted: the framework derives the
observables (masses, couplings, mixings) but carries the Lagrangian form as an external
structure. The hosted SM dynamics (the Lagrangian and its structure) is thus a boundary,
while the derived SM observables (the values read from the D96 moments) are derived physics.
\end{theorem}

\begin{proof}
The SM-dynamics audits \cite{c13:QG242,c13:QG245} establish the classification: the gauge
\emph{symmetry} is derived, but the gauge \emph{dynamics} (the Lagrangian, vertices,
propagators) is hosted; the SM dynamics is not complete as a derivation. The derived
observables --- the fermion sectors, couplings, mixings, weak scale, and precision
predictions (Chapter~\ref{c11:chap:sm}) --- are reads of the D96 moments. Hence the SM
dynamics is a boundary and the observables are derived physics.
\end{proof}

\section{Derived Physics vs Boundary Questions}
\label{c13:sec:derived-vs-boundary}

\begin{theorem}[The classification of remaining items]
\label{c13:thm:classification}
The remaining items of the framework classify into derived physics, experimental
validation, and boundaries: the derived physics is established; the validation items
($S,T,U$, $a_e$, $0\nu\beta\beta$) are derived predictions awaiting experiment; the
boundaries are the ontological, numerical, and hosted limits.
\end{theorem}

\begin{proof}
The post-frontier audit \cite{c13:QG299} classifies the remaining items; VALID001
\cite{c13:VALID001} establishes that $S,T,U$, $a_e$, and $0\nu\beta\beta$ are \emph{derived
predictions} (status A) belonging to the experimental-validation category, with no missing
derivation steps. Hence none is an open derivation issue or a boundary; the boundaries are
those documented in this chapter, and the remaining items partition into derived physics
(established), validation (derived, awaiting experiment), and boundaries (the limits).
\end{proof}

\begin{corollary}[No validation item is an open derivation issue]
\label{c13:cor:no-open-validation}
The quantities $S,T,U$, the electron g-2 $a_e$, and the Majorana character $0\nu\beta\beta$
are not listed as open derivation issues: their derivations are complete, and their
experimental validation is the open task.
\end{corollary}

\section{What Is Not Claimed}
\label{c13:sec:not-claimed}

\begin{theorem}[The framework's non-claims]
\label{c13:thm:non-claims}
The framework does not claim: (i) a derivation of the standard-model Lagrangian (hosted);
(ii) a derivation of the exact Bekenstein $1/4$ coefficient (requires the imported $2\pi$
factor); (iii) a resolution of the ultimate ontological status of $\psi$; (iv) any result
beyond the accepted derivations. Equivalently, the framework claims exactly what it
derives: the physics observables established in the preceding chapters, with the boundaries
disclosed and the validation items awaiting experiment.
\end{theorem}

\begin{proof}
(i) The SM dynamics is hosted \cite{c13:QG242,c13:QG245} (Theorem~\ref{c13:thm:sm-hosted}). (ii) The
Bekenstein $1/4$ coefficient requires the imported $2\pi$ factor \cite{c13:QG185,c13:QG196}
(Theorem~\ref{c13:thm:bekenstein-2pi}). (iii) The ontological status of $\psi$ is unresolved
(Theorem~\ref{c13:thm:psi-face}). (iv) The monograph is assembly-only, restricted to
accepted results. Hence the non-claims are explicit and the claimed content is the derived
physics.
\end{proof}

\section{Consequences for Future Work}
\label{c13:sec:future-work}

\begin{theorem}[The boundaries delimit the future work]
\label{c13:thm:future-work}
The boundaries delimit the future work of the framework: the experimental validation of the
derived predictions, the resolution of the ontological status of $\psi$, and the question
of whether the imported $2\pi$ factor can be internalized. The boundaries are not blockers:
the derived content of the theory is complete within its scope, and the boundaries mark
where the framework honestly stops.
\end{theorem}

\begin{proof}
The derived predictions ($S,T,U$, $a_e$, $0\nu\beta\beta$) await experimental validation
\cite{c13:VALID001}; the ontological status of $\psi$ is a boundary question
(Theorem~\ref{c13:thm:psi-face}); the Bekenstein $2\pi$ factor is an imported boundary
(Theorem~\ref{c13:thm:bekenstein-2pi}). By Lemma~\ref{c13:lem:boundaries-not-failures} the
boundaries are not failures, and the derived physics is complete within the canonical
hierarchy (Chapter~\ref{c12:chap:cosmology}, Corollary~\ref{c12:cor:canonical-cosmology}).
Hence the boundaries are not blockers but the framework's declared edge.
\end{proof}

\section{Summary}
\label{c13:sec:boundaries-summary}

This chapter has documented the boundaries of the theory and shown that they are honest
limits, not derivation failures. Difference is the irreducible primitive boundary; $\psi$ is
the tensor face of Difference with unresolved ontological status; $\pi$ is a numerical
boundary, redundant as a theory input; the Bekenstein $1/4$ coefficient is explicitly tied
to the imported $2\pi$ factor; and the SM Lagrangian is hosted while the observables are
derived. The remaining items partition into derived physics, validation, and boundary, with
the validation items ($S,T,U$, $a_e$, $0\nu\beta\beta$) distinct from open derivation issues
(Theorem~\ref{c13:thm:classification}, Corollary~\ref{c13:cor:no-open-validation}); the
framework claims exactly its derivations (Theorem~\ref{c13:thm:non-claims}); and the
boundaries delimit the honest future work without blocking it
(Theorem~\ref{c13:thm:future-work}).

The boundaries of The Actualization Theory are therefore: Difference (the ontological
boundary and primitive), $\psi$ (the tensor face with unresolved ontological status),
$\pi$ (the numerical boundary), the Bekenstein $2\pi$ factor (the imported quantum
boundary), and the hosted standard-model dynamics --- each an explicit limit, none a
derivation failure. The derived physics is complete within its scope; the experimental
validation of the derived predictions is the open task; and the boundaries mark where the
framework honestly stops. This closes the Boundary Layer; the following chapter documents
the frontier and the falsification paths.

\begin{thebibliography}{99}
\bibitem{c13:QG185} AT-QG Phase 185, \emph{Bekenstein Quarter Coefficient Origin}, Docs/Research.
\bibitem{c13:QG196} AT-QG Phase 196, \emph{Quarter Coefficient Impossibility Proof}, Docs/Research.
\bibitem{c13:QG242} AT-QG Phase 242, \emph{Standard Model Dynamics Audit}, Docs/Research.
\bibitem{c13:QG245} AT-QG Phase 245, \emph{Higgs Yukawa Origin Audit}, Docs/Research.
\bibitem{c13:QG278} AT-QG Phase 278, \emph{Fundamental Boundary Audit}, Docs/Research.
\bibitem{c13:QG286} AT-QG Phase 286, \emph{Difference Duality Audit}, Docs/Research.
\bibitem{c13:QG291} AT-QG Phase 291, \emph{Framework Necessity Audit}, Docs/Research.
\bibitem{c13:QG299} AT-QG Phase 299, \emph{Remaining Frontier Re-audit}, Docs/Research.
\bibitem{c13:MONO007} AT-MONO007, \emph{Referee-Readiness Resolution}, Docs/Zenodo/Monograph.
\bibitem{c13:VALID001} AT-VALID001, \emph{Remaining Untested Items Audit}, Docs/Research.
\end{thebibliography}
 after the
%         preamble that defines the theorem environments (or include via the
%         master file).
% ======================================================================

% ---- Theorem environments (define in the master preamble if not present) ----
% \newtheorem{definition}{Definition}[chapter]
% \newtheorem{lemma}[definition]{Lemma}
% \newtheorem{theorem}[definition]{Theorem}
% \newtheorem{corollary}[definition]{Corollary}
% \newtheorem{remark}[definition]{Remark}

\chapter{Boundaries of The Actualization Theory}
\label{c13:chap:boundaries}

\section{Introduction}
\label{c13:sec:boundaries-introduction}

The preceding chapters derived the physics of The Actualization Theory: the foundation, the
spectrum, quantum mechanics, gravity, the standard model, and cosmology. This chapter is the
Boundary Layer of the monograph: it separates \emph{ontological boundaries} from \emph{physics
derivation}, documents the limits of the framework, and shows that the boundaries are
\emph{not derivation failures} --- they are the explicit limits beyond which the framework
does not claim to derive.

The boundaries documented here are: Difference as an ontological boundary \cite{c13:QG278}; the
status of the tensor face $\psi$ \cite{c13:QG286}; the numerical boundary $\pi$ \cite{c13:QG291};
the Bekenstein $2\pi$ boundary \cite{c13:QG185,c13:QG196}; and the hosted standard-model dynamics
\cite{c13:QG242,c13:QG245}. The chapter draws the distinction between \emph{derived physics} and
\emph{boundary questions} \cite{c13:QG299} and states explicitly what the theory does not claim.

Following VALID001 \cite{c13:VALID001}, \emph{experimental validation is not a boundary}: the
derived predictions $S,T,U$, the electron g-2 $a_e$, and the Majorana character $0\nu\beta\beta$
await experiment and are not listed as derivation gaps. Consistent with MONO007 \cite{c13:MONO007},
the presentation is referee-safe and introduces no new physics and no speculation.

\section{The Nature of Scientific Boundaries}
\label{c13:sec:nature-boundaries}

\begin{lemma}[Scientific boundaries are not failures]
\label{c13:lem:boundaries-not-failures}
A \emph{scientific boundary} is a limit of a framework beyond which the framework does not
claim to derive: either an irreducible primitive (an ontological boundary), a numerical
constant not derivable inside the framework (a numerical boundary), or a hosted structure
(an external component carried without derivation). A documented boundary is not a
derivation failure: it is the explicit, honest limit of the framework, disclosed so that no
claim exceeds what is established.
\end{lemma}

\begin{proof}
A framework is characterized as much by what it does not claim as by what it derives: the
disclosure of a boundary is the statement "this framework does not derive X". The
post-frontier audit \cite{c13:QG299} classifies the remaining items as open, partial,
boundary, or methodology, distinguishing the boundaries from the open derivation items.
Hence a boundary is an honest limit, not a failure of the derivation.
\end{proof}

\section{Difference as an Ontological Boundary}
\label{c13:sec:difference-boundary}

Difference is the ontological boundary of the theory: the irreducible primitive whose
count-conservation law is its definitional identity, and which no deeper notion explains
--- the \emph{primitive} boundary, the irreducible starting point rather than an unexplained
gap. The fundamental-boundary audit \cite{c13:QG278} establishes that every attempt to reduce
it reintroduces it (two distinct things use distinct = difference; a boundary is where
things differ; a transition is a before/after difference); by
Chapter~\ref{c01:chap:difference} (Theorem~\ref{c01:thm:fundamental-boundary}) it is the
fundamental boundary of the theory, and by Lemma~\ref{c13:lem:boundaries-not-failures} an
ontological boundary is an honest primitive limit, not a derivation failure.
\label{c13:thm:difference-ontological}

\section{The Status of \texorpdfstring{$\psi$}{psi}}
\label{c13:sec:psi-status}

\begin{theorem}[\texorpdfstring{$\psi$}{psi} is the tensor face of Difference]
\label{c13:thm:psi-face}
The tensor field $\psi$ is the tensor face of Difference: the traceless part of the
trace/traceless decomposition of the one Difference, read against the tensor reference
$\eta$. Its deepest ontological nature --- whether the tensor content is ultimately
irreducible to the scalar count --- is unresolved: a boundary question, not a derivation gap.
\end{theorem}

\begin{proof}
The Difference duality \cite{c13:QG286} establishes that $\rho$ and $\psi$ are the trace and
traceless faces of the one Difference (Chapter~\ref{c01:chap:difference},
Theorem~\ref{c01:thm:duality}); the tensor face is read against the reference $\eta$
(Chapter~\ref{c02:chap:eta}, Theorem~\ref{c02:thm:tensor-read}). Hence $\psi$ is the tensor
face of Difference --- not an independent primitive. The post-frontier audit
\cite{c13:QG299} classifies the $\psi$ question as a boundary item: the observables of the
tensor face are derived, but its ultimate ontological status is a boundary question.
\end{proof}

\section{The \texorpdfstring{$\pi$}{pi} Boundary}
\label{c13:sec:pi-boundary}

\begin{theorem}[\texorpdfstring{$\pi$}{pi} is redundant as a theory input]
\label{c13:thm:pi-redundant}
\emph{$\pi$} is the universal mathematical constant, the geometric ratio of the circle; in
the canonical architecture it is a \emph{numerical boundary} --- a constant not derivable
inside the framework --- and not a primitive. Correspondingly, $\pi$ is redundant as a
theory input: no derived prediction uses $\pi$ as an input, and where it appears it is an
imported numerical factor.
\end{theorem}

\begin{proof}
The framework-necessity audit \cite{c13:QG291} establishes the classification: every derived
observable --- masses, couplings, mixings, cosmological fractions, spectral indices,
acoustic ratios, and the pre-registered predictions --- is a pure ratio of D96 spectral
constants together with the calibration scale. $\pi$ appears only where an imported quantum
factor is required, most notably in the Bekenstein coefficient; as a constant not derivable
inside the framework it is a numerical boundary, distinct from the primitive reference
$\eta$ (Chapter~\ref{c02:chap:eta}, Theorem~\ref{c02:thm:eta-primitive-pi-boundary}), not a
primitive.
\end{proof}

\section{The Bekenstein \texorpdfstring{$2\pi$}{2pi} Boundary}
\label{c13:sec:bekenstein}

\begin{theorem}[The Bekenstein $1/4$ coefficient requires the imported $2\pi$ factor]
\label{c13:thm:bekenstein-2pi}
The exact Bekenstein coefficient $S = A/4$ is a boundary: the structure $S \propto A$,
$M \propto R$, $T \propto 1/R$ is derived, but the exact $1/4$ coefficient requires the
imported quantum factor $T = \kappa/(2\pi)$, which the framework cannot produce internally.
The boundary is therefore explicitly referenced to the imported quantum factor, not
presented as a derivation of the theory.
\end{theorem}

\begin{proof}
The Bekenstein origin \cite{c13:QG185} derives the structure ($S \propto A$, $M \propto R$,
$T \propto 1/R$); the deficit first-law gives $S = A/(8\pi)$, not $1/4$, and the exact $1/4$
requires the $2\pi$ quantum factor $T = \kappa/(2\pi)$, absent in the framework. The
impossibility proof \cite{c13:QG196} confirms that the exact $1/4$ cannot be derived without
importing $\pi$. Consistent with Chapter~\ref{c10:chap:gravity}
(Theorem~\ref{c10:thm:bekenstein-boundary}), the boundary is referenced, not claimed.
\end{proof}

\section{Hosted Standard Model Dynamics}
\label{c13:sec:hosted-sm}

\begin{theorem}[The SM Lagrangian is hosted, not derived]
\label{c13:thm:sm-hosted}
The standard-model \emph{Lagrangian} and its dynamics are hosted: the framework derives the
observables (masses, couplings, mixings) but carries the Lagrangian form as an external
structure. The hosted SM dynamics (the Lagrangian and its structure) is thus a boundary,
while the derived SM observables (the values read from the D96 moments) are derived physics.
\end{theorem}

\begin{proof}
The SM-dynamics audits \cite{c13:QG242,c13:QG245} establish the classification: the gauge
\emph{symmetry} is derived, but the gauge \emph{dynamics} (the Lagrangian, vertices,
propagators) is hosted; the SM dynamics is not complete as a derivation. The derived
observables --- the fermion sectors, couplings, mixings, weak scale, and precision
predictions (Chapter~\ref{c11:chap:sm}) --- are reads of the D96 moments. Hence the SM
dynamics is a boundary and the observables are derived physics.
\end{proof}

\section{Derived Physics vs Boundary Questions}
\label{c13:sec:derived-vs-boundary}

\begin{theorem}[The classification of remaining items]
\label{c13:thm:classification}
The remaining items of the framework classify into derived physics, experimental
validation, and boundaries: the derived physics is established; the validation items
($S,T,U$, $a_e$, $0\nu\beta\beta$) are derived predictions awaiting experiment; the
boundaries are the ontological, numerical, and hosted limits.
\end{theorem}

\begin{proof}
The post-frontier audit \cite{c13:QG299} classifies the remaining items; VALID001
\cite{c13:VALID001} establishes that $S,T,U$, $a_e$, and $0\nu\beta\beta$ are \emph{derived
predictions} (status A) belonging to the experimental-validation category, with no missing
derivation steps. Hence none is an open derivation issue or a boundary; the boundaries are
those documented in this chapter, and the remaining items partition into derived physics
(established), validation (derived, awaiting experiment), and boundaries (the limits).
\end{proof}

\begin{corollary}[No validation item is an open derivation issue]
\label{c13:cor:no-open-validation}
The quantities $S,T,U$, the electron g-2 $a_e$, and the Majorana character $0\nu\beta\beta$
are not listed as open derivation issues: their derivations are complete, and their
experimental validation is the open task.
\end{corollary}

\section{What Is Not Claimed}
\label{c13:sec:not-claimed}

\begin{theorem}[The framework's non-claims]
\label{c13:thm:non-claims}
The framework does not claim: (i) a derivation of the standard-model Lagrangian (hosted);
(ii) a derivation of the exact Bekenstein $1/4$ coefficient (requires the imported $2\pi$
factor); (iii) a resolution of the ultimate ontological status of $\psi$; (iv) any result
beyond the accepted derivations. Equivalently, the framework claims exactly what it
derives: the physics observables established in the preceding chapters, with the boundaries
disclosed and the validation items awaiting experiment.
\end{theorem}

\begin{proof}
(i) The SM dynamics is hosted \cite{c13:QG242,c13:QG245} (Theorem~\ref{c13:thm:sm-hosted}). (ii) The
Bekenstein $1/4$ coefficient requires the imported $2\pi$ factor \cite{c13:QG185,c13:QG196}
(Theorem~\ref{c13:thm:bekenstein-2pi}). (iii) The ontological status of $\psi$ is unresolved
(Theorem~\ref{c13:thm:psi-face}). (iv) The monograph is assembly-only, restricted to
accepted results. Hence the non-claims are explicit and the claimed content is the derived
physics.
\end{proof}

\section{Consequences for Future Work}
\label{c13:sec:future-work}

\begin{theorem}[The boundaries delimit the future work]
\label{c13:thm:future-work}
The boundaries delimit the future work of the framework: the experimental validation of the
derived predictions, the resolution of the ontological status of $\psi$, and the question
of whether the imported $2\pi$ factor can be internalized. The boundaries are not blockers:
the derived content of the theory is complete within its scope, and the boundaries mark
where the framework honestly stops.
\end{theorem}

\begin{proof}
The derived predictions ($S,T,U$, $a_e$, $0\nu\beta\beta$) await experimental validation
\cite{c13:VALID001}; the ontological status of $\psi$ is a boundary question
(Theorem~\ref{c13:thm:psi-face}); the Bekenstein $2\pi$ factor is an imported boundary
(Theorem~\ref{c13:thm:bekenstein-2pi}). By Lemma~\ref{c13:lem:boundaries-not-failures} the
boundaries are not failures, and the derived physics is complete within the canonical
hierarchy (Chapter~\ref{c12:chap:cosmology}, Corollary~\ref{c12:cor:canonical-cosmology}).
Hence the boundaries are not blockers but the framework's declared edge.
\end{proof}

\section{Summary}
\label{c13:sec:boundaries-summary}

This chapter has documented the boundaries of the theory and shown that they are honest
limits, not derivation failures. Difference is the irreducible primitive boundary; $\psi$ is
the tensor face of Difference with unresolved ontological status; $\pi$ is a numerical
boundary, redundant as a theory input; the Bekenstein $1/4$ coefficient is explicitly tied
to the imported $2\pi$ factor; and the SM Lagrangian is hosted while the observables are
derived. The remaining items partition into derived physics, validation, and boundary, with
the validation items ($S,T,U$, $a_e$, $0\nu\beta\beta$) distinct from open derivation issues
(Theorem~\ref{c13:thm:classification}, Corollary~\ref{c13:cor:no-open-validation}); the
framework claims exactly its derivations (Theorem~\ref{c13:thm:non-claims}); and the
boundaries delimit the honest future work without blocking it
(Theorem~\ref{c13:thm:future-work}).

The boundaries of The Actualization Theory are therefore: Difference (the ontological
boundary and primitive), $\psi$ (the tensor face with unresolved ontological status),
$\pi$ (the numerical boundary), the Bekenstein $2\pi$ factor (the imported quantum
boundary), and the hosted standard-model dynamics --- each an explicit limit, none a
derivation failure. The derived physics is complete within its scope; the experimental
validation of the derived predictions is the open task; and the boundaries mark where the
framework honestly stops. This closes the Boundary Layer; the following chapter documents
the frontier and the falsification paths.

\begin{thebibliography}{99}
\bibitem{c13:QG185} AT-QG Phase 185, \emph{Bekenstein Quarter Coefficient Origin}, Docs/Research.
\bibitem{c13:QG196} AT-QG Phase 196, \emph{Quarter Coefficient Impossibility Proof}, Docs/Research.
\bibitem{c13:QG242} AT-QG Phase 242, \emph{Standard Model Dynamics Audit}, Docs/Research.
\bibitem{c13:QG245} AT-QG Phase 245, \emph{Higgs Yukawa Origin Audit}, Docs/Research.
\bibitem{c13:QG278} AT-QG Phase 278, \emph{Fundamental Boundary Audit}, Docs/Research.
\bibitem{c13:QG286} AT-QG Phase 286, \emph{Difference Duality Audit}, Docs/Research.
\bibitem{c13:QG291} AT-QG Phase 291, \emph{Framework Necessity Audit}, Docs/Research.
\bibitem{c13:QG299} AT-QG Phase 299, \emph{Remaining Frontier Re-audit}, Docs/Research.
\bibitem{c13:MONO007} AT-MONO007, \emph{Referee-Readiness Resolution}, Docs/Zenodo/Monograph.
\bibitem{c13:VALID001} AT-VALID001, \emph{Remaining Untested Items Audit}, Docs/Research.
\end{thebibliography}
 after the
%         preamble that defines the theorem environments (or include via the
%         master file).
% ======================================================================

% ---- Theorem environments (define in the master preamble if not present) ----
% \newtheorem{definition}{Definition}[chapter]
% \newtheorem{lemma}[definition]{Lemma}
% \newtheorem{theorem}[definition]{Theorem}
% \newtheorem{corollary}[definition]{Corollary}
% \newtheorem{remark}[definition]{Remark}

\chapter{Boundaries of The Actualization Theory}
\label{c13:chap:boundaries}

\section{Introduction}
\label{c13:sec:boundaries-introduction}

The preceding chapters derived the physics of The Actualization Theory: the foundation, the
spectrum, quantum mechanics, gravity, the standard model, and cosmology. This chapter is the
Boundary Layer of the monograph: it separates \emph{ontological boundaries} from \emph{physics
derivation}, documents the limits of the framework, and shows that the boundaries are
\emph{not derivation failures} --- they are the explicit limits beyond which the framework
does not claim to derive.

The boundaries documented here are: Difference as an ontological boundary \cite{c13:QG278}; the
status of the tensor face $\psi$ \cite{c13:QG286}; the numerical boundary $\pi$ \cite{c13:QG291};
the Bekenstein $2\pi$ boundary \cite{c13:QG185,c13:QG196}; and the hosted standard-model dynamics
\cite{c13:QG242,c13:QG245}. The chapter draws the distinction between \emph{derived physics} and
\emph{boundary questions} \cite{c13:QG299} and states explicitly what the theory does not claim.

Following VALID001 \cite{c13:VALID001}, \emph{experimental validation is not a boundary}: the
derived predictions $S,T,U$, the electron g-2 $a_e$, and the Majorana character $0\nu\beta\beta$
await experiment and are not listed as derivation gaps. Consistent with MONO007 \cite{c13:MONO007},
the presentation is referee-safe and introduces no new physics and no speculation.

\section{The Nature of Scientific Boundaries}
\label{c13:sec:nature-boundaries}

\begin{lemma}[Scientific boundaries are not failures]
\label{c13:lem:boundaries-not-failures}
A \emph{scientific boundary} is a limit of a framework beyond which the framework does not
claim to derive: either an irreducible primitive (an ontological boundary), a numerical
constant not derivable inside the framework (a numerical boundary), or a hosted structure
(an external component carried without derivation). A documented boundary is not a
derivation failure: it is the explicit, honest limit of the framework, disclosed so that no
claim exceeds what is established.
\end{lemma}

\begin{proof}
A framework is characterized as much by what it does not claim as by what it derives: the
disclosure of a boundary is the statement "this framework does not derive X". The
post-frontier audit \cite{c13:QG299} classifies the remaining items as open, partial,
boundary, or methodology, distinguishing the boundaries from the open derivation items.
Hence a boundary is an honest limit, not a failure of the derivation.
\end{proof}

\section{Difference as an Ontological Boundary}
\label{c13:sec:difference-boundary}

Difference is the ontological boundary of the theory: the irreducible primitive whose
count-conservation law is its definitional identity, and which no deeper notion explains
--- the \emph{primitive} boundary, the irreducible starting point rather than an unexplained
gap. The fundamental-boundary audit \cite{c13:QG278} establishes that every attempt to reduce
it reintroduces it (two distinct things use distinct = difference; a boundary is where
things differ; a transition is a before/after difference); by
Chapter~\ref{c01:chap:difference} (Theorem~\ref{c01:thm:fundamental-boundary}) it is the
fundamental boundary of the theory, and by Lemma~\ref{c13:lem:boundaries-not-failures} an
ontological boundary is an honest primitive limit, not a derivation failure.
\label{c13:thm:difference-ontological}

\section{The Status of \texorpdfstring{$\psi$}{psi}}
\label{c13:sec:psi-status}

\begin{theorem}[\texorpdfstring{$\psi$}{psi} is the tensor face of Difference]
\label{c13:thm:psi-face}
The tensor field $\psi$ is the tensor face of Difference: the traceless part of the
trace/traceless decomposition of the one Difference, read against the tensor reference
$\eta$. Its deepest ontological nature --- whether the tensor content is ultimately
irreducible to the scalar count --- is unresolved: a boundary question, not a derivation gap.
\end{theorem}

\begin{proof}
The Difference duality \cite{c13:QG286} establishes that $\rho$ and $\psi$ are the trace and
traceless faces of the one Difference (Chapter~\ref{c01:chap:difference},
Theorem~\ref{c01:thm:duality}); the tensor face is read against the reference $\eta$
(Chapter~\ref{c02:chap:eta}, Theorem~\ref{c02:thm:tensor-read}). Hence $\psi$ is the tensor
face of Difference --- not an independent primitive. The post-frontier audit
\cite{c13:QG299} classifies the $\psi$ question as a boundary item: the observables of the
tensor face are derived, but its ultimate ontological status is a boundary question.
\end{proof}

\section{The \texorpdfstring{$\pi$}{pi} Boundary}
\label{c13:sec:pi-boundary}

\begin{theorem}[\texorpdfstring{$\pi$}{pi} is redundant as a theory input]
\label{c13:thm:pi-redundant}
\emph{$\pi$} is the universal mathematical constant, the geometric ratio of the circle; in
the canonical architecture it is a \emph{numerical boundary} --- a constant not derivable
inside the framework --- and not a primitive. Correspondingly, $\pi$ is redundant as a
theory input: no derived prediction uses $\pi$ as an input, and where it appears it is an
imported numerical factor.
\end{theorem}

\begin{proof}
The framework-necessity audit \cite{c13:QG291} establishes the classification: every derived
observable --- masses, couplings, mixings, cosmological fractions, spectral indices,
acoustic ratios, and the pre-registered predictions --- is a pure ratio of D96 spectral
constants together with the calibration scale. $\pi$ appears only where an imported quantum
factor is required, most notably in the Bekenstein coefficient; as a constant not derivable
inside the framework it is a numerical boundary, distinct from the primitive reference
$\eta$ (Chapter~\ref{c02:chap:eta}, Theorem~\ref{c02:thm:eta-primitive-pi-boundary}), not a
primitive.
\end{proof}

\section{The Bekenstein \texorpdfstring{$2\pi$}{2pi} Boundary}
\label{c13:sec:bekenstein}

\begin{theorem}[The Bekenstein $1/4$ coefficient requires the imported $2\pi$ factor]
\label{c13:thm:bekenstein-2pi}
The exact Bekenstein coefficient $S = A/4$ is a boundary: the structure $S \propto A$,
$M \propto R$, $T \propto 1/R$ is derived, but the exact $1/4$ coefficient requires the
imported quantum factor $T = \kappa/(2\pi)$, which the framework cannot produce internally.
The boundary is therefore explicitly referenced to the imported quantum factor, not
presented as a derivation of the theory.
\end{theorem}

\begin{proof}
The Bekenstein origin \cite{c13:QG185} derives the structure ($S \propto A$, $M \propto R$,
$T \propto 1/R$); the deficit first-law gives $S = A/(8\pi)$, not $1/4$, and the exact $1/4$
requires the $2\pi$ quantum factor $T = \kappa/(2\pi)$, absent in the framework. The
impossibility proof \cite{c13:QG196} confirms that the exact $1/4$ cannot be derived without
importing $\pi$. Consistent with Chapter~\ref{c10:chap:gravity}
(Theorem~\ref{c10:thm:bekenstein-boundary}), the boundary is referenced, not claimed.
\end{proof}

\section{Hosted Standard Model Dynamics}
\label{c13:sec:hosted-sm}

\begin{theorem}[The SM Lagrangian is hosted, not derived]
\label{c13:thm:sm-hosted}
The standard-model \emph{Lagrangian} and its dynamics are hosted: the framework derives the
observables (masses, couplings, mixings) but carries the Lagrangian form as an external
structure. The hosted SM dynamics (the Lagrangian and its structure) is thus a boundary,
while the derived SM observables (the values read from the D96 moments) are derived physics.
\end{theorem}

\begin{proof}
The SM-dynamics audits \cite{c13:QG242,c13:QG245} establish the classification: the gauge
\emph{symmetry} is derived, but the gauge \emph{dynamics} (the Lagrangian, vertices,
propagators) is hosted; the SM dynamics is not complete as a derivation. The derived
observables --- the fermion sectors, couplings, mixings, weak scale, and precision
predictions (Chapter~\ref{c11:chap:sm}) --- are reads of the D96 moments. Hence the SM
dynamics is a boundary and the observables are derived physics.
\end{proof}

\section{Derived Physics vs Boundary Questions}
\label{c13:sec:derived-vs-boundary}

\begin{theorem}[The classification of remaining items]
\label{c13:thm:classification}
The remaining items of the framework classify into derived physics, experimental
validation, and boundaries: the derived physics is established; the validation items
($S,T,U$, $a_e$, $0\nu\beta\beta$) are derived predictions awaiting experiment; the
boundaries are the ontological, numerical, and hosted limits.
\end{theorem}

\begin{proof}
The post-frontier audit \cite{c13:QG299} classifies the remaining items; VALID001
\cite{c13:VALID001} establishes that $S,T,U$, $a_e$, and $0\nu\beta\beta$ are \emph{derived
predictions} (status A) belonging to the experimental-validation category, with no missing
derivation steps. Hence none is an open derivation issue or a boundary; the boundaries are
those documented in this chapter, and the remaining items partition into derived physics
(established), validation (derived, awaiting experiment), and boundaries (the limits).
\end{proof}

\begin{corollary}[No validation item is an open derivation issue]
\label{c13:cor:no-open-validation}
The quantities $S,T,U$, the electron g-2 $a_e$, and the Majorana character $0\nu\beta\beta$
are not listed as open derivation issues: their derivations are complete, and their
experimental validation is the open task.
\end{corollary}

\section{What Is Not Claimed}
\label{c13:sec:not-claimed}

\begin{theorem}[The framework's non-claims]
\label{c13:thm:non-claims}
The framework does not claim: (i) a derivation of the standard-model Lagrangian (hosted);
(ii) a derivation of the exact Bekenstein $1/4$ coefficient (requires the imported $2\pi$
factor); (iii) a resolution of the ultimate ontological status of $\psi$; (iv) any result
beyond the accepted derivations. Equivalently, the framework claims exactly what it
derives: the physics observables established in the preceding chapters, with the boundaries
disclosed and the validation items awaiting experiment.
\end{theorem}

\begin{proof}
(i) The SM dynamics is hosted \cite{c13:QG242,c13:QG245} (Theorem~\ref{c13:thm:sm-hosted}). (ii) The
Bekenstein $1/4$ coefficient requires the imported $2\pi$ factor \cite{c13:QG185,c13:QG196}
(Theorem~\ref{c13:thm:bekenstein-2pi}). (iii) The ontological status of $\psi$ is unresolved
(Theorem~\ref{c13:thm:psi-face}). (iv) The monograph is assembly-only, restricted to
accepted results. Hence the non-claims are explicit and the claimed content is the derived
physics.
\end{proof}

\section{Consequences for Future Work}
\label{c13:sec:future-work}

\begin{theorem}[The boundaries delimit the future work]
\label{c13:thm:future-work}
The boundaries delimit the future work of the framework: the experimental validation of the
derived predictions, the resolution of the ontological status of $\psi$, and the question
of whether the imported $2\pi$ factor can be internalized. The boundaries are not blockers:
the derived content of the theory is complete within its scope, and the boundaries mark
where the framework honestly stops.
\end{theorem}

\begin{proof}
The derived predictions ($S,T,U$, $a_e$, $0\nu\beta\beta$) await experimental validation
\cite{c13:VALID001}; the ontological status of $\psi$ is a boundary question
(Theorem~\ref{c13:thm:psi-face}); the Bekenstein $2\pi$ factor is an imported boundary
(Theorem~\ref{c13:thm:bekenstein-2pi}). By Lemma~\ref{c13:lem:boundaries-not-failures} the
boundaries are not failures, and the derived physics is complete within the canonical
hierarchy (Chapter~\ref{c12:chap:cosmology}, Corollary~\ref{c12:cor:canonical-cosmology}).
Hence the boundaries are not blockers but the framework's declared edge.
\end{proof}

\section{Summary}
\label{c13:sec:boundaries-summary}

This chapter has documented the boundaries of the theory and shown that they are honest
limits, not derivation failures. Difference is the irreducible primitive boundary; $\psi$ is
the tensor face of Difference with unresolved ontological status; $\pi$ is a numerical
boundary, redundant as a theory input; the Bekenstein $1/4$ coefficient is explicitly tied
to the imported $2\pi$ factor; and the SM Lagrangian is hosted while the observables are
derived. The remaining items partition into derived physics, validation, and boundary, with
the validation items ($S,T,U$, $a_e$, $0\nu\beta\beta$) distinct from open derivation issues
(Theorem~\ref{c13:thm:classification}, Corollary~\ref{c13:cor:no-open-validation}); the
framework claims exactly its derivations (Theorem~\ref{c13:thm:non-claims}); and the
boundaries delimit the honest future work without blocking it
(Theorem~\ref{c13:thm:future-work}).

The boundaries of The Actualization Theory are therefore: Difference (the ontological
boundary and primitive), $\psi$ (the tensor face with unresolved ontological status),
$\pi$ (the numerical boundary), the Bekenstein $2\pi$ factor (the imported quantum
boundary), and the hosted standard-model dynamics --- each an explicit limit, none a
derivation failure. The derived physics is complete within its scope; the experimental
validation of the derived predictions is the open task; and the boundaries mark where the
framework honestly stops. This closes the Boundary Layer; the following chapter documents
the frontier and the falsification paths.

\begin{thebibliography}{99}
\bibitem{c13:QG185} AT-QG Phase 185, \emph{Bekenstein Quarter Coefficient Origin}, Docs/Research.
\bibitem{c13:QG196} AT-QG Phase 196, \emph{Quarter Coefficient Impossibility Proof}, Docs/Research.
\bibitem{c13:QG242} AT-QG Phase 242, \emph{Standard Model Dynamics Audit}, Docs/Research.
\bibitem{c13:QG245} AT-QG Phase 245, \emph{Higgs Yukawa Origin Audit}, Docs/Research.
\bibitem{c13:QG278} AT-QG Phase 278, \emph{Fundamental Boundary Audit}, Docs/Research.
\bibitem{c13:QG286} AT-QG Phase 286, \emph{Difference Duality Audit}, Docs/Research.
\bibitem{c13:QG291} AT-QG Phase 291, \emph{Framework Necessity Audit}, Docs/Research.
\bibitem{c13:QG299} AT-QG Phase 299, \emph{Remaining Frontier Re-audit}, Docs/Research.
\bibitem{c13:MONO007} AT-MONO007, \emph{Referee-Readiness Resolution}, Docs/Zenodo/Monograph.
\bibitem{c13:VALID001} AT-VALID001, \emph{Remaining Untested Items Audit}, Docs/Research.
\end{thebibliography}
